\section{Five Dynasties regional and frontier ledger events}
\label{sec:five-dynasties-ten-kingdoms-event-anchors-batch-two}
这些锚点对应账本中可复核的日期与来源链。正史、文书、碑志和考古材料各有观察范围；宏只提供正文导航，不将政策、战报或册命直接换算为全体社会事实。

\subsection{梁、吴与两浙并行的二十年，908--930年}

唐亡后的第二年，开封的后梁仍须用汴州的军粮、运河和亲军维持中原朝廷；淮南的吴、两浙的吴越、巴蜀的前蜀与岐地的李茂贞，却以各自的都城、将校和年号安排同一时期。908--912年的册命、称王和继承，不能被读成后梁向四方逐次扩张的年表。江淮的水路、两浙的海港与巴蜀的山川，分别限定了征粮、用兵和信息往来的速度；一项来自开封的诏令只能说明中央的主张，不能自动说明地方如何执行。%
\footnote{后梁初年、吴越与吴的政权活动可参《旧五代史》梁纪、《新五代史》梁本纪及吴、吴越世家、《资治通鉴》后梁纪，并与《吴越备史》、十国国别史辑佚材料互校。中原史书的正统框架与国别史的成书、传本均须分别辨认。}

916年契丹建立政权，918年前后北方军镇和江南诸国的交涉更加复杂；北方的战争并不停止南方的继承与筑城。923年后唐入汴灭梁，925年又灭前蜀，表面上拉近了中原和巴蜀的政治距离；926年庄宗死于兵变，旋即证明城池易手、官府接收和军费供给不能被一场胜利同时解决。唐庄宗的伶官故事有强烈的后世道德化形状，更可靠的提问是军赏、军纪和洛阳、汴州之间的政治资源为何未能稳定结合。%
\footnote{契丹建国、后唐灭梁、灭前蜀及926年兵变见《旧五代史》唐书、《新五代史》唐本纪与契丹传、《资治通鉴》后梁、后唐纪。辽方材料与中原材料对称谓、日期和政治目的的差异，不应由一方叙事静默抹平。}

927--930年间，后唐中枢、荆南、楚、闽、吴与吴越继续在不同的税源、航道和家族网络中运作。它们不是中原王朝空隙里的静态“十国背景”：江陵控制长江中游的转运，潭州连接湖湘资源，福州和广州面向海上网络，金陵则依托江淮交通。国别史散佚、墓志与钱币材料分布不均，使这些地区比开封更难写成逐年齐整的叙事；这种材料不均衡应被保留，而不能由后来的统一结局填平。%
\footnote{十国并行政治的年序以《新五代史》十国世家、《资治通鉴》为检索骨架，并须参照《吴越备史》、南唐史料、碑志、钱币与城市考古。单一墓志、窖藏或港口遗址只能约束特定地点和时段，不能代表整个政权的财政或社会。}

\eventanchor{V03-908-ANNAL-CENTRAL}{V03-908-ANNAL-CENTRAL} 
\eventanchor{V03-908-CROSS-POLITY}{V03-908-CROSS-POLITY} 
\eventanchor{V03-908-INSTITUTION-SOCIETY}{V03-908-INSTITUTION-SOCIETY} 
\eventanchor{V03-909-ANNAL-CENTRAL}{V03-909-ANNAL-CENTRAL} 
\eventanchor{V03-909-CROSS-POLITY}{V03-909-CROSS-POLITY} 
\eventanchor{V03-909-INSTITUTION-SOCIETY}{V03-909-INSTITUTION-SOCIETY} 
\eventanchor{V03-910-ANNAL-CENTRAL}{V03-910-ANNAL-CENTRAL} 
\eventanchor{V03-910-CROSS-POLITY}{V03-910-CROSS-POLITY} 
\eventanchor{V03-910-INSTITUTION-SOCIETY}{V03-910-INSTITUTION-SOCIETY} 
\eventanchor{V03-911-ANNAL-CENTRAL}{V03-911-ANNAL-CENTRAL} 
\eventanchor{V03-911-CROSS-POLITY}{V03-911-CROSS-POLITY} 
\eventanchor{V03-911-INSTITUTION-SOCIETY}{V03-911-INSTITUTION-SOCIETY} 
\eventanchor{V03-912-ANNAL-CENTRAL}{V03-912-ANNAL-CENTRAL} 
\eventanchor{V03-912-CROSS-POLITY}{V03-912-CROSS-POLITY} 
\eventanchor{V03-912-INSTITUTION-SOCIETY}{V03-912-INSTITUTION-SOCIETY} 
\eventanchor{V03-913-ANNAL-CENTRAL}{V03-913-ANNAL-CENTRAL} 
\eventanchor{V03-913-CROSS-POLITY}{V03-913-CROSS-POLITY} 
\eventanchor{V03-913-INSTITUTION-SOCIETY}{V03-913-INSTITUTION-SOCIETY} 
\eventanchor{V03-914-ANNAL-CENTRAL}{V03-914-ANNAL-CENTRAL} 
\eventanchor{V03-914-CROSS-POLITY}{V03-914-CROSS-POLITY} 
\eventanchor{V03-914-INSTITUTION-SOCIETY}{V03-914-INSTITUTION-SOCIETY} 
\eventanchor{V03-915-ANNAL-CENTRAL}{V03-915-ANNAL-CENTRAL} 
\eventanchor{V03-915-CROSS-POLITY}{V03-915-CROSS-POLITY} 
\eventanchor{V03-915-INSTITUTION-SOCIETY}{V03-915-INSTITUTION-SOCIETY} 
\eventanchor{V03-916-ANNAL-CENTRAL}{V03-916-ANNAL-CENTRAL} 
\eventanchor{V03-916-CROSS-POLITY}{V03-916-CROSS-POLITY} 
\eventanchor{V03-916-INSTITUTION-SOCIETY}{V03-916-INSTITUTION-SOCIETY} 
\eventanchor{V03-917-ANNAL-CENTRAL}{V03-917-ANNAL-CENTRAL} 
\eventanchor{V03-917-CROSS-POLITY}{V03-917-CROSS-POLITY} 
\eventanchor{V03-917-INSTITUTION-SOCIETY}{V03-917-INSTITUTION-SOCIETY} 
\eventanchor{V03-918-ANNAL-CENTRAL}{V03-918-ANNAL-CENTRAL} 
\eventanchor{V03-918-CROSS-POLITY}{V03-918-CROSS-POLITY} 
\eventanchor{V03-918-INSTITUTION-SOCIETY}{V03-918-INSTITUTION-SOCIETY} 
\eventanchor{V03-919-ANNAL-CENTRAL}{V03-919-ANNAL-CENTRAL} 
\eventanchor{V03-919-CROSS-POLITY}{V03-919-CROSS-POLITY} 
\eventanchor{V03-919-INSTITUTION-SOCIETY}{V03-919-INSTITUTION-SOCIETY} 
\eventanchor{V03-920-ANNAL-CENTRAL}{V03-920-ANNAL-CENTRAL} 
\eventanchor{V03-920-CROSS-POLITY}{V03-920-CROSS-POLITY} 
\eventanchor{V03-920-INSTITUTION-SOCIETY}{V03-920-INSTITUTION-SOCIETY} 
\eventanchor{V03-921-ANNAL-CENTRAL}{V03-921-ANNAL-CENTRAL} 
\eventanchor{V03-921-CROSS-POLITY}{V03-921-CROSS-POLITY} 
\eventanchor{V03-921-INSTITUTION-SOCIETY}{V03-921-INSTITUTION-SOCIETY} 
\eventanchor{V03-922-ANNAL-CENTRAL}{V03-922-ANNAL-CENTRAL} 
\eventanchor{V03-922-CROSS-POLITY}{V03-922-CROSS-POLITY} 
\eventanchor{V03-922-INSTITUTION-SOCIETY}{V03-922-INSTITUTION-SOCIETY} 
\eventanchor{V03-923-ANNAL-CENTRAL}{V03-923-ANNAL-CENTRAL} 
\eventanchor{V03-923-CROSS-POLITY}{V03-923-CROSS-POLITY} 
\eventanchor{V03-923-INSTITUTION-SOCIETY}{V03-923-INSTITUTION-SOCIETY} 
\eventanchor{V03-924-ANNAL-CENTRAL}{V03-924-ANNAL-CENTRAL} 
\eventanchor{V03-924-CROSS-POLITY}{V03-924-CROSS-POLITY} 
\eventanchor{V03-924-INSTITUTION-SOCIETY}{V03-924-INSTITUTION-SOCIETY} 
\eventanchor{V03-925-ANNAL-CENTRAL}{V03-925-ANNAL-CENTRAL} 
\eventanchor{V03-925-CROSS-POLITY}{V03-925-CROSS-POLITY} 
\eventanchor{V03-925-INSTITUTION-SOCIETY}{V03-925-INSTITUTION-SOCIETY} 
\eventanchor{V03-926-ANNAL-CENTRAL}{V03-926-ANNAL-CENTRAL} 
\eventanchor{V03-926-CROSS-POLITY}{V03-926-CROSS-POLITY} 
\eventanchor{V03-926-INSTITUTION-SOCIETY}{V03-926-INSTITUTION-SOCIETY} 
\eventanchor{V03-927-ANNAL-CENTRAL}{V03-927-ANNAL-CENTRAL} 
\eventanchor{V03-927-CROSS-POLITY}{V03-927-CROSS-POLITY} 
\eventanchor{V03-927-INSTITUTION-SOCIETY}{V03-927-INSTITUTION-SOCIETY} 
\eventanchor{V03-928-ANNAL-CENTRAL}{V03-928-ANNAL-CENTRAL} 
\eventanchor{V03-928-CROSS-POLITY}{V03-928-CROSS-POLITY} 
\eventanchor{V03-928-INSTITUTION-SOCIETY}{V03-928-INSTITUTION-SOCIETY} 
\eventanchor{V03-929-ANNAL-CENTRAL}{V03-929-ANNAL-CENTRAL} 
\eventanchor{V03-929-CROSS-POLITY}{V03-929-CROSS-POLITY} 
\eventanchor{V03-929-INSTITUTION-SOCIETY}{V03-929-INSTITUTION-SOCIETY} 
\eventanchor{V03-930-ANNAL-CENTRAL}{V03-930-ANNAL-CENTRAL} 
\eventanchor{V03-930-CROSS-POLITY}{V03-930-CROSS-POLITY} 
